\section{Method}
\label{sec:method}

We now describe the estimator used throughout the remainder of the thesis.
The construction proceeds in two stages: first we define a spatially
weighted intensity estimate that corrects for the irregular sensor density
illustrated in Figure~\ref{fig:network}, and second we derive an interval
estimator for the recurrence time from that intensity, together with a
confidence procedure that accounts for the finite number of observed
events.

\subsection{Spatially Weighted Intensity Estimation}
Let $s_1, \dots, s_n$ denote the locations of the $n$ deployed sensors, and
let $E_i$ denote the set of event times recorded by sensor $i$ during the
observation window $[0, T]$. Under the independence assumption discussed in
Section~\ref{sec:background}, each sensor observes a thinned version of the
global event process, with thinning probability proportional to a local
coverage function $\pi(s_i)$. We estimate the aggregate intensity as a
coverage-corrected sum over sensors,
\begin{equation}
\label{eq:intensity}
\hat{\lambda}(t) = \sum_{i=1}^{n} \frac{1}{\pi(s_i)} \sum_{\tau \in E_i}
K_h(t - \tau),
\end{equation}
where $K_h$ is a smoothing kernel with bandwidth $h$ and $\pi(s_i)$ is
estimated from the sensor placement study described in
Section~\ref{sec:background}, following the placement-aware correction of
\citet{fictional2020}. The bandwidth $h$ controls the bias--variance
trade-off of the estimate in the usual way: small values of $h$ track rapid
changes in the underlying rate but produce a noisy estimate, whereas large
values smooth over genuine rate variation.

Several practical considerations affect the choice of kernel and bandwidth
in this setting:
\begin{itemize}
  \item The bandwidth should be chosen relative to the shortest recurrence
  interval of interest, since a bandwidth much larger than the interval of
  interest will average away the signal entirely.
  \item Boundary effects near $t = 0$ and $t = T$ bias the raw kernel
  estimate downward; we correct for this using the standard reflection
  method.
  \item When sensor dropout is correlated with local event intensity, the
  coverage correction in Equation~\eqref{eq:intensity} is no longer
  unbiased, and a more elaborate missing-data model is required.
\end{itemize}

\subsection{Recurrence Interval Estimation}
Given the smoothed intensity estimate $\hat{\lambda}(t)$, we define the
estimated recurrence interval at time $t$ as the reciprocal of the local
intensity, $\hat{r}(t) = 1 / \hat{\lambda}(t)$. A first-order delta-method
argument yields an approximate confidence interval for $\hat{r}(t)$ in terms
of the variance of $\hat{\lambda}(t)$, which we estimate using a block
bootstrap over sensor-level event sequences to account for spatial
correlation between nearby sensors. This bootstrap procedure resamples
whole sensors rather than individual events, which preserves the
dependence structure within each sensor's event sequence and yields
substantially more conservative, and we believe more realistic, confidence
intervals than a naive event-level bootstrap would produce.

\subsection{Evaluation on Synthetic Data}
We first validate the estimator on synthetic event sequences generated from
a known Hawkes process with a time-varying baseline intensity, so that the
true recurrence interval is known exactly and can be compared against the
estimate produced by Equation~\eqref{eq:intensity}. Across a range of
sensor densities and dropout rates, the coverage-corrected estimator
recovers the true recurrence interval with markedly lower bias than an
uncorrected estimator that ignores the spatial coverage function, confirming
that the correction term derived from \citet{fictional2020} is necessary
rather than a merely theoretical refinement.

\subsection{Field Study}
We then apply the estimator to a field-collected dataset gathered from the
sensor deployment shown in Figure~\ref{fig:network}, spanning several
months of continuous monitoring. The estimated recurrence intervals from
the field data are broadly consistent with maintenance records collected
independently by facility operators, which we take as external validation
of the method beyond the synthetic study. Discrepancies that do arise are
concentrated near sensors with the lowest local coverage, consistent with
the variance inflation predicted by the block-bootstrap confidence
procedure described above, and consistent with the general caution urged
by \citep{daley2003} regarding inference from sparsely sampled point
processes. The remaining chapters of the thesis build on this estimator to
address the downstream problem of scheduling preventive maintenance under
uncertainty in the estimated recurrence interval.
